\documentclass[11pt]{amsart}

\usepackage[T1]{fontenc}
\usepackage{lmodern}
\usepackage{microtype}
\usepackage{mathtools,amssymb,amsthm}
\usepackage[hidelinks]{hyperref}

\hypersetup{
  pdftitle={A 3x3 Counterexample Family to Conjecture 4.5 of Chen-Cao's arXiv:1801.00225v1}
}

\newtheorem{theorem}{Theorem}
\newtheorem{lemma}[theorem]{Lemma}
\theoremstyle{remark}
\newtheorem*{remark}{Remark}

\DeclareMathOperator{\per}{per}
\newcommand{\ws}{\omega}
\newcommand{\wt}{\tilde{\omega}}

\title[A counterexample family to Conjecture 4.5 of \texttt{arXiv:1801.00225v1}]
{A $3\times3$ Counterexample Family to Conjecture~4.5 of
Chen--Cao's \texttt{arXiv:1801.00225v1}}

\date{}

\subjclass[2020]{15A15, 15B51, 05A05, 52B12}
\keywords{permanent, doubly substochastic matrix, prescribed matrix sum,
counterexample, convexity on a polytope}

\begin{document}

\begin{abstract}
Let $\ws_n$ be the set of $n\times n$ doubly substochastic matrices (nonnegative,
all row sums and all column sums at most $1$), let $\sigma(A)$ be the sum of all
entries of $A$, and let $\ws_n^{\,s}=\{A\in\ws_n:\sigma(A)=s\}$. Chen and Cao
determined $\max\{\per(I-A):A\in\ws_n^{\,s}\}$ for $n$ even and for $n$ odd with
$s\le n-1$, left the remaining cell open, and conjectured a two-branch value for
$n=3$, $2<s\le 3$ (their Conjecture~4.5), whose printed form needs two readings fixed
before it can be tested (Section~\ref{sec:src}). In the fixed-mass reading, with the
second branch read as the one their own Lemma~4.4 proves, the conjecture is false at
\emph{every} $s$ in the open interval $(2,3)$. The witness is the one-parameter
family
\[
 A(s)=\begin{pmatrix}0&1-u&u\\ 1-u&0&u\\ u&u&0\end{pmatrix},
 \qquad u=\tfrac{s-2}{2},
\]
\[
 \per\bigl(I-A(s)\bigr)=\frac{s^{3}-5s^{2}+4s+12}{4},
\]
which lies in $\ws_3^{\,s}$ and exceeds \emph{both} of the conjectured branches, the
excess over the first being exactly $s(s-3)^{2}/4$. At $s=5/2$ the conjecture
predicts $23/16$ and the witness gives $51/32$. Equivalently, Chen--Cao's own
Lemma~4.4 is not tight anywhere on $(2,3)$. The true value
of $\max\{\per(I-A):A\in\ws_3^{\,s}\}$ remains open; at $s=5/2$ what is established
here is $51/32\le\max\le2$. The journal version of record was not obtained, so every
statement here about what the source prints is a statement about the e-print
\texttt{arXiv:1801.00225v1}.
\end{abstract}

\maketitle

\section{The statement refuted}\label{sec:src}

\subsection*{Notation, in the source's own words}
An $n\times n$ nonnegative matrix is \emph{doubly stochastic} if every row sum and
every column sum equals $1$, and \emph{doubly substochastic} if every row sum and
every column sum is at most~$1$. Write $\Omega_n$ for the doubly stochastic
matrices, $\ws_n$ for the doubly substochastic ones, and $\wt_n$ for the
\emph{row} substochastic ones (only the row sums are constrained). For a square
matrix $A=[a_{ij}]$ let $\sigma(A)=\sum_{i,j}a_{ij}$, and put
\[
 \ws_n^{\,s}=\{A\in\ws_n:\sigma(A)=s\},
 \qquad
 \wt_n^{\,s}=\{A\in\wt_n:\sigma(A)=s\}.
\]
No symmetry, no zero diagonal and no rationality is assumed anywhere in $\ws_n$.
The distinction between $\ws_n$ and $\wt_n$ is the axis of this note: Chen and Cao
settled the row substochastic problem in the cell at issue and left the doubly
substochastic one open.

\subsection*{The source statement, verbatim}
Chen and Cao~\cite{chencao} prove
\begin{equation}\label{eq:thm31}
 \max\{\per(I-A):A\in\ws_n^{\,s}\}
 =2^{e/2}\Bigl[1+\bigl(\tfrac{s-e}{2}\bigr)^{2}\Bigr],
 \qquad e=2\lfloor s/2\rfloor,
\end{equation}
under the hypothesis ``$n$ is even, or $n$ is odd and $\sigma(A)\le n-1$'' (their
Theorem~3.1; their Theorem~3.13 is the same formula over $\wt_n^{\,s}$), and they
say of the excluded cell that ``for $A\in\ws_n^s$, the upper bound of $\per(I-A)$
still remains mystery''. In Section~4 they settle the row substochastic half of
that cell,
\begin{equation}\label{eq:thm41}
 \max\{\per(I-A):A\in\wt_n^{\,s},\ n-1<s\le n\}=2^{\frac{n-1}{2}}
 \qquad (n \text{ odd}),
\end{equation}
observe that their maximiser there is \emph{not} doubly substochastic (a column sum
exceeds $1$), and then conjecture the $n=3$ value.

All statement numbers and all quotations below are from the \LaTeX{} e-print
\texttt{arXiv:1801.00225v1} (\url{https://arxiv.org/e-print/1801.00225}), whose single
source file \texttt{per\_I-A\_LAASubmission.tex} has SHA-256 beginning
\texttt{9113694030113dda}; the statement numbering follows from its preamble, in which
\texttt{thm}, \texttt{lem} and \texttt{conj} share one counter. The target statement,
verbatim:
\begin{equation}\label{eq:c45}
 \max_{A\in\ws_3}\{\per(I-A)\}=
 \begin{cases}
  \dfrac{\sigma^{2}(A)-5\sigma(A)+12}{4}
   &\text{if }\ \dfrac{-3+\sqrt{57}}{2}<\sigma(A)\le3,\\[8pt]
  6-4\sigma(A)
   &\text{if }\ 2<\sigma(A)\le\dfrac{-3+\sqrt{57}}{2},
 \end{cases}
\end{equation}
and the half of it the authors \emph{prove}, verbatim (Lemma~4.4):
\begin{equation}\label{eq:l44}
 \max_{A\in\ws_3}\{\per(I-A)\}
 \ \ge\ \max\Bigl\{\tfrac{\sigma^{2}(A)-5\sigma(A)+12}{4},\ 6-2\sigma(A)\Bigr\}
 \qquad(2<\sigma(A)\le3),
\end{equation}
proved by exhibiting the two matrices
\begin{equation}\label{eq:their}
 A_0=\begin{pmatrix}0&\tfrac12&\tfrac{\sigma}{2}-1\\
   \tfrac12&0&\tfrac12\\ \tfrac{\sigma}{2}-1&\tfrac12&0\end{pmatrix},
 \qquad
 A_1=\begin{pmatrix}0&1&0\\1&0&0\\0&0&\sigma-2\end{pmatrix},
\end{equation}
with $\per(I-A_0)=\frac{\sigma^{2}-5\sigma+12}{4}$ and
$\per(I-A_1)=6-2\sigma$. \textbf{So \eqref{eq:c45} is precisely the assertion that
\eqref{eq:l44} is tight}, and that is what we refute.

\subsection*{Two readings, and which one is refuted}
Two points of reading must be fixed before anything is claimed.

First, \eqref{eq:c45} is printed as a maximum over all of $\ws_3$ whose right-hand
side depends on $\sigma(A)$; read literally that is ill-formed, since
$\max_{A\in\ws_3}\per(I-A)$ is the single number~$2$. The intended
\emph{fixed-mass} reading, $\max$ over $\ws_3^{\,\sigma}$, is forced by
\eqref{eq:l44} directly above it and by Conjecture~4.6, which writes
$\max\{\per(I-B):B\in\ws_3^{\,s'}\}$ explicitly. That is the reading refuted here.

Second, the printed second branch $6-4\sigma$ is negative for $\sigma>3/2$ and hence
throughout its own stated range, whereas $\per(I-A)\ge\det(I-A)\ge0$ for doubly
substochastic $A$; it is also contradicted in print by \eqref{eq:l44}, which the
source states just above it. It is a misprint for the $6-2\sigma$ of \eqref{eq:l44}, and three
independent arguments say so: the printed threshold $\frac{-3+\sqrt{57}}2$ is the root
of $s^{2}+3s-12$, i.e.\ the crossover of the first branch with $6-2s$ (the crossover
with $6-4s$ is the root of $s^{2}+11s-12$, namely $s=1$, outside $(2,3]$
entirely); the authors' own $A_1$ realises $6-2\sigma$ exactly and $6-4\sigma$ never;
and positivity forbids $6-4\sigma$. \textbf{Refuting the printed reading would
therefore be worth nothing, and we do not claim it as content.} We refute the
corrected, and strongest, reading
\begin{equation}\label{eq:V}
 V(s):=\max\Bigl\{\tfrac{s^{2}-5s+12}{4},\ 6-2s\Bigr\},
\end{equation}
which is exactly the lower bound \eqref{eq:l44} proves. As it happens the witness
below exceeds $6-4s$ as well, so the refutation is independent of how the misprint
is read; and at the headline mass $s=5/2$ the governing branch is the first one,
printed identically in the \emph{proved} \eqref{eq:l44}, so there the two readings
coincide.

\subsection*{What text we read}
\textbf{The version of record was not obtained.} Every sentence of this note about
``what the source prints'' is a statement about \texttt{arXiv:1801.00225v1} and
nothing else. The arXiv abstract page carries \texttt{[v1]} only and no
\texttt{journal-ref}; we found no open copy of the journal version; and the reference
list deposited with \texttt{doi:10.1016/j.laa.2018.06.031} agrees with the preprint's
twelve \verb|\bibitem|s in eleven entries and differs in one, a term-rank paper
replacing a preprint citation, so refereeing did alter the document. The Crossref
record for the DOI carries no update or relation entry and we found no corrigendum
notice; we do not conclude from that that no correction exists, since the journal text
itself was not read and one citation-index channel could not be read either. We
therefore state the claim in the form: \emph{Conjecture~4.5 of
\texttt{arXiv:1801.00225v1}, in the reading fixed above, is false.} The residual risk
is discussed in Section~\ref{sec:open}.

\section{The witness}\label{sec:wit}

For $s\in(2,3)$ set $u=\frac{s-2}{2}\in(0,\tfrac12)$ and
\begin{equation}\label{eq:A}
 A(s)=\begin{pmatrix}0&1-u&u\\ 1-u&0&u\\ u&u&0\end{pmatrix}.
\end{equation}
Membership in $\ws_3^{\,s}$ is checkable by eye. The entries are $0$, $u>0$ and
$1-u>0$. The row sums are $(1,\,1,\,2u)$ and, since $A(s)$ is symmetric, the column
sums are the same; $2u=s-2\le1$ exactly because $s\le3$. Finally
$\sigma(A(s))=2(1-u)+4u=2+2u=s$, exactly. The column-sum clause is the one that
matters: dropping it would place the matrix in the row substochastic class
$\wt_3$, where the maximum \eqref{eq:thm41} is already known.

For a zero-diagonal $A=\begin{psmallmatrix}0&a&b\\c&0&d\\e&f&0\end{psmallmatrix}$ the
six permutations of $S_3$ give
\begin{equation}\label{eq:per3}
 \per(I-A)=1+ac+be+df-ade-bcf ,
\end{equation}
so with $a=c=1-u$ and $b=d=e=f=u$,
\begin{equation}\label{eq:f}
\begin{split}
 \per\bigl(I-A(s)\bigr)&=1+(1-u)^{2}+2u^{2}-2(1-u)u^{2}=2-2u+u^{2}+2u^{3}\\
 &=:f(s)=\frac{s^{3}-5s^{2}+4s+12}{4}.
\end{split}
\end{equation}

\subsection*{The headline instance}
At $s=5/2$, i.e.\ $u=1/4$,
\[
 A=\begin{pmatrix}0&\tfrac34&\tfrac14\\[2pt] \tfrac34&0&\tfrac14\\[2pt]
   \tfrac14&\tfrac14&0\end{pmatrix}\in\ws_3,
 \qquad
 \sigma(A)=\tfrac52,
\]
with row sums and column sums $(1,1,\tfrac12)$ and
\[
 \per(I-A)=1+\tfrac9{16}+\tfrac1{16}+\tfrac1{16}-\tfrac3{64}-\tfrac3{64}
 =\tfrac{51}{32}=1.59375 .
\]
Since $\tfrac52>\tfrac{-3+\sqrt{57}}2=2.274917\ldots$, the first branch of
\eqref{eq:c45} is in force and predicts $\tfrac{25/4-25/2+12}{4}=\tfrac{23}{16}
=1.4375$. The excess is exactly $\tfrac5{32}$ and the ratio is $\tfrac{51}{46}$.

\subsection*{Exact values}
Table~\ref{tab:inst} lists every value we record. All entries are exact rationals.

\begin{table}[ht]
\small
\begin{tabular}{llllll}
\hline
$s$ & $1-u,\ u$ & $\per(I-A(s))=f(s)$ & $V(s)$ & branch & excess\\
\hline
$21/10$ & $19/20,\ 1/20$ & $7611/4000=1.902750$ & $9/5=1.8$ & second & $411/4000$\\
$9/4$   & $7/8,\ 1/8$    & $453/256=1.769531$   & $3/2=1.5$ & second & $69/256$\\
$5/2$   & $3/4,\ 1/4$    & $51/32=1.593750$     & $23/16=1.4375$ & first & $5/32$\\
$11/4$  & $5/8,\ 3/8$    & $383/256=1.496094$   & $93/64=1.453125$ & first & $11/256$\\
$14/5$  & $3/5,\ 2/5$    & $186/125=1.488000$   & $73/50=1.46$ & first & $7/250$\\
$29/10$ & $11/20,\ 9/20$ & $5939/4000=1.484750$ & $591/400=1.4775$ & first & $29/4000$\\
$2$ & $1,\ 0$ & $2$ & $2$ & --- & $0$\\
$3$ & $1/2,\ 1/2$ & $3/2$ & $3/2$ & --- & $0$\\
\hline
\end{tabular}
\medskip
\caption{The witness against the corrected conjectured value \eqref{eq:V}. In the last
two rows the mass is an endpoint, where the true maximum is published and the family
returns exactly it, with no violation: $f(2)=2$ is \eqref{eq:thm31} at $n=3$, $e=2$,
and $f(3)=3/2$ is Chen--Cao's Lemma~4.2, the value over $\Omega_3$.}
\label{tab:inst}
\end{table}

\section{The refutation}\label{sec:thm}

\begin{theorem}\label{thm:main}
For every $s$ with $2<s<3$ the matrix $A(s)$ of \eqref{eq:A} lies in
$\ws_3^{\,s}$ and
\[
 \max\{\per(I-A):A\in\ws_3^{\,s}\}\ \ge\ f(s)\ >\ V(s),
\]
where $V$ is the value \eqref{eq:V}: Conjecture~4.5 of \texttt{arXiv:1801.00225v1} in
the reading fixed in Section~\ref{sec:src}, which is also the lower bound Lemma~4.4
proves.
Consequently Conjecture~4.5 of \texttt{arXiv:1801.00225v1}, in the reading fixed in
Section~\ref{sec:src}, is false at every mass in $(2,3)$, and Lemma~4.4 is not tight
anywhere on that interval.
\end{theorem}

\begin{proof}
Membership and the value $f(s)$ are Section~\ref{sec:wit}. Two polynomial identities
finish it; both are exact coefficient identities, not numerical comparisons. First,
\begin{equation}\label{eq:id1}
 f(s)-\frac{s^{2}-5s+12}{4}=\frac{s^{3}-6s^{2}+9s}{4}=\frac{s(s-3)^{2}}{4}>0
 \qquad(2<s<3),
\end{equation}
since $s>0$ and $s\ne3$. Second, with $t=s-2\in(0,1)$,
\begin{equation}\label{eq:id2}
 f(s)-(6-2s)=\frac{s^{3}-5s^{2}+12s-12}{4}=t+\frac{t^{2}}4+\frac{t^{3}}4>0 .
\end{equation}
Because $V$ is the \emph{maximum} of the two branches and $f$ exceeds both of them at
every interior mass, $f>V$ there. (For completeness, $f$ also exceeds the printed
branch $6-4s$, which is negative on the whole interval, so no reading of the misprint
and no reading of the threshold rescues the conjecture.) Equality in \eqref{eq:id1}
occurs on $[2,3]$ only at $s=3$, and \eqref{eq:id2} degenerates only at $s=2$.
\end{proof}

The witness respects the two upper bounds proved for this functional: since
$A(s)\in\ws_3\subset\wt_3$, Chen--Cao's \eqref{eq:thm41} and Malek's
theorem~\cite{malek} both give $\per(I-A(s))\le2$ at $n=3$, and indeed
$2-f(s)=(s-2)(2+3s-s^{2})/4>0$ on $(2,3]$, the roots of $2+3s-s^{2}$ being
$\tfrac{3\pm\sqrt{17}}2$.

\section{What is not settled}\label{sec:open}

\subsection*{The upper bound is open, and this is the main limit}
What is proved here is $\max\ge f(s)$ together with $f(s)>V(s)$. It is \emph{not}
proved that $\max=f(s)$. Malek's bound~\cite{malek} gives only $\max\le2$ at $n=3$,
and $f<2$ on $(2,3]$, so it does not decide the question. \textbf{At $s=5/2$ the
honest bracket is $51/32\le\max\le2$}; the conjectured $23/16$ is excluded, and
nothing here determines the true value.

\subsection*{No optimality of $A(s)$ is claimed}
$A(s)$ is exhibited only as a feasible point beating $V(s)$. \textbf{It is not claimed
here to maximise $\per(I-A)$ over $\ws_3^{\,s}$}, on any slice of that polytope or
over all of it; no global search over the nine free entries is offered here, and none
is performed by the program accompanying this note. The natural tool for a rigorous
finish is Cao and Chen's partition of the polytope of doubly substochastic matrices by
the sum of all entries, which characterises the extreme points of exactly
$\ws_n^{\,s}$~\cite{caochen}. That $A(s)$ has a zero diagonal is a fact about it;
nothing here assumes a zero diagonal, and Lemma~4.4's own second witness $A_1$ of
\eqref{eq:their} carries $a_{33}=\sigma-2\ne0$.

\subsection*{The residual risk in the journal text}
Only \texttt{arXiv:1801.00225v1} was read, and the reference list deposited with the
DOI shows that it is not the printed document (Section~\ref{sec:src}). The refutation
is unaffected by an edit to the misprinted branch or to the
threshold, since at $s=5/2$ the printed and the corrected readings coincide and the
governing first branch there is printed identically in the \emph{proved}
\eqref{eq:l44}. The risk that remains is wholesale replacement of the conjecture in
print: if refereeing replaced the two-branch value by the cubic $f$, then this note
shrinks to a correction of the preprint record. The refutation of the v1 text is
unaffected either way.

\subsection*{Prior art}
The source has four indexed citers, all of which were read; none of them quotes
Conjecture~4.5 or~4.6 at all, they cite Lemma~4.2 and Theorems~3.1, 3.13 and~4.1.
Malek's theorem~\cite{malek}, which we read only through its zbMATH review, lands
adjacent --- an upper bound of $2^{\lfloor n/2\rfloor}$, with no prescribed-mass
parameter --- rather than on the target. The citer list is not certified complete: one
citation-index channel could not be read, and the zbMATH review of the source's own
record is licence-blocked; neither is read here as a negative.

The closest live work is Lavi~\cite{lavi}, which settles a \emph{neighbouring}
conjecture of the same source --- the Kim--Roush value
$\max_{A\in\Omega_{2k+1}}\per(I-A)=3\cdot2^{k-2}$, i.e.\ the doubly stochastic
endpoint $s=n$ --- with an equality classification. It contains no counterexample to
Conjecture~4.5 and no value on $\ws_3^{\,s}$ for $2<s<3$; it records the cell at issue
here as open: ``For odd $n$ and $n-1<s\le n$, they determined the row-substochastic
maximum but left the doubly substochastic case open.'' Two further papers by
overlapping authors are adjacent but do not touch the target: Chen, Li, Yang, Zhu and
Cao~\cite{clyzc19} prove convexity of the permanent along the segment $[I,A]$ for
\emph{row} substochastic $A$ and conjecture a minimum over transportation polytopes,
with no maximum over $\ws_n^{\,s}$; and Chen, Zhu, Li, Yang and Cao~\cite{czlyc21}
characterise the extreme points of transportation polytopes with fixed total sum and
give the maximum of $\per(A)$ there, not of $\per(I-A)$.

\section{Reproduction and verification}\label{sec:repro}

The load-bearing verification is by hand. Section~\ref{sec:wit} exhibits a
$3\times3$ matrix with entries in $\{0,\tfrac14,\tfrac34\}$, six row and column sums,
and a six-term permanent; Section~\ref{sec:thm} is two polynomial identities in one
variable. No computation is needed to check any claim in this note, and none of the
claims rests on floating-point arithmetic.

A program accompanies this note as a re-derivation, not as evidence: it recomputes every
permanent by the full $n!$ expansion over exact rationals, reproduces
Table~\ref{tab:inst} including the governing branch at each mass, and verifies
\eqref{eq:id1} and \eqref{eq:id2} by exact coefficient comparison rather than by
sampling. It also carries out checks on material not stated in this note --- the
geometry of the symmetric zero-diagonal slice at each tabulated mass, and a direct-sum
lift of the witness to $n=5,7$ --- which support no claim made above. Its transcript
records $97$ passing checks and no failures.

What that run does and does not establish should be stated plainly. Its per-matrix
checks --- membership in $\ws_3$, the mass and the permanent --- are carried out at the
eight masses of Table~\ref{tab:inst}, six interior and two endpoints, and at no other
mass; the strict comparison $f(s)>V(s)$ is checked at the six interior masses only,
while at $s=2$ and $s=3$ the program checks and records the equality $f(s)=V(s)$, with
excess exactly $0$ and no violation. Several of its check names and step headings are
phrased as statements over the whole open interval, but what is evaluated at those
places is exactly those sampled masses. The claim for every $s\in(2,3)$ rests on the two
polynomial identities of Section~\ref{sec:thm}, which are proved by hand and re-checked
by the program as identities between coefficients rather than at sample points; the
strict signs of those two expressions on $(2,3)$ are read off by hand and are not
machine-verified. The program does \emph{not} attempt the open upper bound of
Section~\ref{sec:open}, performs no global search, contains no floating-point decision,
and reads no source text: what the source prints is quoted from
\texttt{arXiv:1801.00225v1} in Section~\ref{sec:src}, and the journal version of record
was never obtained.

\begin{thebibliography}{9}

\bibitem{chencao}
Z.~Chen and L.~Cao,
\emph{On the maximum of the permanent of $(I-A)$},
Linear Algebra Appl. \textbf{555} (2018), 412--431.
\href{https://doi.org/10.1016/j.laa.2018.06.031}{doi:10.1016/j.laa.2018.06.031};
\href{https://arxiv.org/abs/1801.00225}{arXiv:1801.00225}.
All statement numbers and quotations in this note are from the e-print
\texttt{arXiv:1801.00225v1}, identified in Section~\ref{sec:src}; it is
the only full text we could reach, and the journal version was not obtained.

\bibitem{malek}
Malek,
\emph{On the maximum of $\mathrm{per}(I-A)$},
Linear Multilinear Algebra \textbf{19} (1986), no.~4, 347--355.
Read only through its zbMATH review, Zbl 0604.15009 (reviewer F.~J. Gaines), from
which we quote the bound $\per(I-A)\le2^{\lfloor n/2\rfloor}$ for row substochastic
$A$; the article itself is paywalled and pre-digital and we did not obtain it.

\bibitem{caochen}
L.~Cao and Z.~Chen,
\emph{Partitions of the polytope of doubly substochastic matrices},
Linear Algebra Appl. (2019).
\href{https://doi.org/10.1016/j.laa.2018.10.024}{doi:10.1016/j.laa.2018.10.024};
\href{https://arxiv.org/abs/1803.00435}{arXiv:1803.00435}.
Cited for its characterisation of the extreme points of $\ws_n^{\,s}$, which is the
natural tool for the upper bound left open in Section~\ref{sec:open}.

\bibitem{lavi}
Y.~Lavi,
\emph{The maximum of $\per(I-A)$ in odd order},
\href{https://arxiv.org/abs/2608.08933}{arXiv:2608.08933} (2026).

\bibitem{clyzc19}
Z.~Chen, J.~Li, L.~Yang, Z.~Zhu and L.~Cao,
\emph{Inequalities for permanents and permanental minors of row substochastic
matrices},
Electron. J. Linear Algebra \textbf{35} (2019), 633--643.
\href{https://doi.org/10.13001/ela.2019.5143}{doi:10.13001/ela.2019.5143}.

\bibitem{czlyc21}
Z.~Chen, Z.~Zhu, J.~Li, L.~Yang and L.~Cao,
\emph{Extreme points of certain transportation polytopes with fixed total sums},
Electron. J. Linear Algebra \textbf{37} (2021), 256--271.
\href{https://doi.org/10.13001/ela.2021.5141}{doi:10.13001/ela.2021.5141}.

\end{thebibliography}

\end{document}
